\documentclass[../../main.tex]{subfiles}
\begin{document}
\section{Family 3: Bochner, heat flow, and the \texorpdfstring{$H^{-1}$}{H-1} calculus}
\label{sec:family-bochner}

\paragraph{Object followed.}
The first eigenspace of the diffusion generator, and the energy of derivatives of test functions,
measured in the negative Sobolev norm attached to that generator.

\paragraph{What it buys.}
Two things that no other family supplies. First, the conversion device: an inequality that turns
``curvature $t$ plus covariance $\norm{\Cov\mu}_\op$'' into a spectral gap, sharper than
Bakry--\'Emery. This is the box marked \emph{Improved Lichnerowicz} in
Figure~\ref{fig:kls-architecture} and it is the reason short-time localization is usable at all.
Second, a calculus in which coordinate functions and quadratic functions have computable spectral
mass, which is what makes the moment-map input of Section~\ref{sec:family-moment-map} consumable.

\subsection{The improved Lichnerowicz inequality}
\label{subsec:improved-lichnerowicz}

Let $\mu=e^{-U}\dd x$ with $\Hess U\succeq tI$. Bakry--\'Emery, equivalently Brascamp--Lieb, gives
$\CP(\mu)\le1/t$. Klartag's sharper geometric-mean estimate is the following
\cite{Klartag2023Logarithmic}.

\begin{theorem}[Improved Lichnerowicz; Klartag]
\label{thm:improved-lichnerowicz}
Let $\mu=e^{-U}\dd x$ be a probability measure on $\R^n$ with $\Hess U\succeq tI$ for some $t>0$.
Then
\begin{equation}
  \CP(\mu)\le\sqrt{\frac{\norm{\Cov\mu}_\op}{t}} .
  \label{eq:improved-lichnerowicz}
\end{equation}
\end{theorem}

This is imported, not proved here; the argument is short enough to record, because the shape of it
explains where each factor in the bridge \eqref{eq:kls-bridge} comes from.

\begin{proof}[Sketch, following {\cite{Klartag2023Logarithmic}}]
Let $-Lf=\lambda f$ with $f$ a normalized first nonconstant eigenfunction, so that
$\int|\nabla f|^2\dd\mu=\lambda$ and $\lambda=\CP(\mu)^{-1}$. The integrated Bochner formula,
combined with the Poincar\'e inequality applied to each partial derivative of $f$, gives
\begin{equation}
  \int\inner{\Hess U\,\nabla f}{\nabla f}\dd\mu
  \le\lambda\Bigl|\int\nabla f\dd\mu\Bigr|^2 .
  \label{eq:bochner-step}
\end{equation}
Integration by parts bounds the right-hand side by the covariance:
\begin{equation}
  \Bigl|\int\nabla f\dd\mu\Bigr|^2\le\lambda^2\norm{\Cov\mu}_\op .
  \label{eq:ibp-step}
\end{equation}
Strong convexity bounds the left-hand side from below by $t\int|\nabla f|^2\dd\mu=t\lambda$.
Chaining the three gives $t\lambda\le\lambda^3\norm{\Cov\mu}_\op$, that is
$\lambda^2\ge t/\norm{\Cov\mu}_\op$, which is \eqref{eq:improved-lichnerowicz}.
\end{proof}

Note the two places the argument spends: \eqref{eq:bochner-step} is where a \emph{gradient} energy
is converted into information about $\int\nabla f\dd\mu$, a single vector, and
\eqref{eq:ibp-step} is where that vector is charged to the operator norm of the covariance. Both
steps are sharp for the Gaussian; neither knows anything about $f$ beyond its first moments. That
is the sense in which this family ``averages away'' the test function, and it is the source of the
limitation recorded at the end of this section.

\subsection{The \texorpdfstring{$H^{-1}$}{H-1} norm and the Barthe--Klartag inequality}
\label{subsec:hminus1}

Let $L=\Delta-\nabla V\cdot\nabla$ be the reversible generator of $\mu=e^{-V}\dd x$, so that
$\CP(\mu)=\lambda_1(-L)^{-1}$. For centered $g$ define
\begin{equation}
  \norm{g}_{H^{-1}(\mu)}^2
  =\inner{g}{(-L)^{-1}g}
  =\int_0^\infty\inner{e^{sL}g}{g}\dd s
  =\int_{\lambda_1}^\infty\frac{\dd\nu_g(\lambda)}\lambda ,
  \label{eq:hminus1-def}
\end{equation}
where $\nu_g$ is the spectral measure of $g$. The last expression is the useful one: the
$H^{-1}$ norm is the spectral mass of $g$ weighted by $1/\lambda$, so it is large exactly when $g$
loads the bottom of the spectrum.

Barthe and Klartag proved that, when the derivatives of $f$ are centered
\cite{BartheKlartag2019SpectralGaps},
\begin{equation}
  \Var_\mu f\le\sum_i\norm{\partial_if}_{H^{-1}(\mu)}^2 .
  \label{eq:barthe-klartag}
\end{equation}
Applied to $f(x)=|x|^2$, whose derivatives are the coordinate functions up to a factor $2$, this
gives
\begin{equation}
  \Var(|X|^2)\le4\sum_i\norm{x_i}_{H^{-1}(\mu)}^2 ,
  \label{eq:bk-radial}
\end{equation}
so the thin-shell problem becomes a statement about the low spectral mass of \emph{coordinate}
functions. This is the entry point through which the moment-map estimates of
Section~\ref{sec:family-moment-map} are consumed.

\begin{remark}[Why quadratics are special]
\label{rem:quadratics-special}
The derivatives of a quadratic function are linear. Inequality \eqref{eq:barthe-klartag} therefore
reduces a quadratic test function to $H^{-1}$ norms of \emph{linear} functions, which the Stein
kernel controls directly (Section~\ref{subsec:mm-stein}). For a general $f$ the derivatives
$\partial_if$ are arbitrary functions and \eqref{eq:barthe-klartag} reduces nothing. This one
sentence is the quadratic-to-all-functions gap, in its sharpest available form.
\end{remark}

\subsection{Heat flow and spectral monotonicity}
\label{subsec:heat-flow}

Smooth the measure by Gaussian convolution, $\mu_s=\mu*\gamma_s$, and define the conditional
expectation operator
\begin{equation}
  Q_sf(y)=\frac{P_s(f\rho)(y)}{P_s\rho(y)}=\E\bigl[f(X)\mid X+\sqrt sG=y\bigr],
  \label{eq:Qs-def}
\end{equation}
with $\rho$ the density of $\mu$ and $P_s$ the heat semigroup. The operator $Q_s$ is the
deterministic, time-reversed counterpart of stochastic localization; concretely,
\begin{equation}
  \sum_i\norm{Q_sx_i}_{L^2(\mu_s)}^2=\E|a_{1/s}|^2,
  \qquad
  \frac{\dd}{\dd t}\E|a_t|^2=\E\Tr(A_t^2),
  \label{eq:Qs-localization-dictionary}
\end{equation}
the second identity being \eqref{eq:sl-centroid-sde} in expectation. So the two families are two
readings of the same object.

Klartag and Putterman proved that the Rayleigh quotient of $Q_sg$ decreases under the heat flow
\cite{KlartagPutterman2021SpectralMonotonicity}; consequently low-frequency modes cannot disappear
too fast,
\begin{equation}
  \norm{Q_sg}_2^2\ge e^{-sE(g)}\norm{g}_2^2 ,
  \label{eq:spectral-monotonicity}
\end{equation}
with $E(g)$ the Rayleigh quotient of $g$. This converts covariance estimates into bounds on the
low spectral mass of coordinate functions; integrating that mass against $1/\lambda$ as in
\eqref{eq:hminus1-def} is how Klartag--Lehec obtained their polylogarithmic thin-shell and KLS
bounds \cite{KlartagLehec2022Polylog}.

\begin{question}[An unresolved operator comparison]
\label{q:literature-PsQs}
Klartag--Lehec prove Rayleigh-quotient and spectral-trace inequalities related to, but weaker than,
the operator inequality
\[
  P_sQ_s\succeq e^{sL} .
\]
A suitable strengthening would remove some of the spectral-projection losses in their argument.
This is a question from the literature \cite{KlartagLehec2022Polylog}, recorded here for
orientation; it is not a repository target and carries no ledger node.
\end{question}

\paragraph{The precise missing estimate.}
Uniform control of $\norm{\partial_if}_{H^{-1}}$ for the derivatives of an \emph{arbitrary} test
function, rather than for coordinates and for derivatives of quadratics.

\paragraph{Why it stalls.}
Every step above is an averaging step. Bochner in
\eqref{eq:bochner-step} keeps only $\int\nabla f\dd\mu$; \eqref{eq:barthe-klartag} keeps only the
spectral masses of the $\partial_if$; \eqref{eq:spectral-monotonicity} is a statement about a
single Rayleigh quotient. Trace, coordinate, or averaged spectral information does not bound every
slow mode, and a near-extremizing eigenfunction of a general log-concave measure is exactly the
object about which these averages say least. Route~S (Group~4) is the attempt to keep the
eigenfunction itself in the argument rather than averaging it away.
\end{document}
